% HAND-WRITTEN fragment -- placed by Adc.tex. Numbers from tab_adc_power.tex
% and tab_adc_corners.tex.
%
% REWRITTEN 2026-08-25 together with tab_adc_power.tex; see that file's header
% for why the previous numbers were withdrawn.  Source:
% ~/chips/castalia/simruns/adcchar_20260825/README.md S2 (corner power) and S3
% (the two protocols), with S5 (the maestro_char headless run, 9 points, which
% reproduces the as-delivered row to 0.02 uA).

\subsubsection{Supply Current and Conversion Energy}\label{sss:adc-power}

Power depends on which clock protocol the converter is driven with, so both are
reported (Table~\ref{tab:adc-power}). On the stimulus as delivered a conversion
takes \SI{4.600}{\micro\second} and costs \SI{12.95}{\micro\watt},
\SI{59.6}{\pico\joule} per conversion; at the design rate the same converter
draws \SI{18.09}{\micro\watt} and \SI{20.8}{\pico\joule}. Four times the
throughput for \SI{40}{\percent} more power is the expected trade --- the static
term is carried four times longer in the slow case --- and the energy per
conversion is the figure that changes, by \(2.9\times\).

Power corners by \(2.7\times\) where the transfer does not
(Table~\ref{tab:adc-corners}): \SIrange{15.8}{43.1}{\micro\watt} across the five
process points, the slow corner alone a factor of two above the other four. That
signature is leakage and bias rather than switching, consistent with the
\SI{1.0}{\volt} core logic slowing down.

Currents are averaged over one whole conversion, and they depend on the sweep
step: the nominal point reads \SI{6.36}{\micro\ampere} on a 63-code step and
\SI{12.56}{\micro\ampere} on a 220-code step, because the switching term
dominates. A converter current quoted without its code step is not a comparable
number.
